%=================================%
\section{Conclusions}
\label{sec:conclusion}
%=================================%

In this paper we have conducted a systematic parameter sweep of 2D, locally isothermal eccentric planet-disc interactions. We have directly solved the global linearised wave equations for a variety of disc temperature profiles, surface densities and aspect ratios using the methodology described in \cite{FairbairnRafikov_2022}. Our direct numerical calculation of modes circumvents the approximations associated with the classical WKB calculations. The resulting wave structure allows us to compute the net torque, which we further decompose into the Lindblad and linear corotation contributions. We perform this for planetary eccentricities ranging from 0.0 to 0.12. In particular, this allows us to probe planetary orbital evolution as the planet begins to move supersonically relative to the gas. The ratio $\tilde{e} = e/h_\textrm{p}$ is the critical measure of this transition and hence our run with the lowest value of $h_\textrm{p} = 0.06$ allows us to interrogate this transition best. 

As $\tilde{e}$ exceeds unity, the torque density undergoes clear changes in its radial profile, with the inner and outer disc contributions reversing in sign. This is notably accompanied by a net torque reversal where angular momentum is now injected into the planetary orbit. However, a net torque increase does not necessarily yield outwards migration: within the range of parameter space investigated, the outwards migration only happens for the highest eccentricities in thin (low $h_\mathrm{p}$) discs with shallow surface density gradients (small $p$). More generally, in the transonic regime (at $\tilde e\sim 1$) we see a strong enhancement in eccentricity and semi-major axis damping rates. The rapid eccentricity damping in the transonic case (roughly $10^2$ times faster than the semi-major axis damping) should efficiently circularize high-$e$ planets and eventually force them to migrate inwards. 

We have also bench-marked our approach by performing an extensive sweep through the space of relevant disc parameters, the results of which were compared with previous analytical and numerical investigations of eccentric planet-disc interaction. Our results appear to be consistent with the existing numerical calculations, although future 2D and 3D simulations across a range of disc models would be desirable to test this further. This will also help address the caveats intrinsic to our 2D, linear methodology which includes the uncertain softening parameter as well as the neglect of nonlinear horseshoe corotation contributions.

Finally, we make our full results freely available in the supplementary materials for future testing and for implementation by the community in e.g. the exoplanetary population synthesis studies.


